\documentclass[10pt,a4paper]{amsart}
\usepackage[margin=19mm]{geometry}
\usepackage{amsmath,amssymb,mathtools}
\usepackage[scr=rsfs,cal=euler]{mathalfa}
\usepackage{xfrac}
\usepackage[unicode,hidelinks,bookmarks=false]{hyperref}
\newcommand{\ie}{i.e.\ }
\newcommand{\cf}{cf.\ }
\newcommand{\resp}{respectively\ }
\newcommand{\Subsection}{\subsection}
\providecommand{\nobreakdash}{\nobreak\hbox{-}}
\setlength{\emergencystretch}{2em}
\multlinegap0pt
\usepackage{bm}
\usepackage{bbm}
\usepackage{dsfont}
\usepackage{xspace}
\usepackage{cancel}
\usepackage{mathtools}
\usepackage{amsthm}
\makeatletter
\@ifundefined{theorem}{\newtheorem{theorem}{Theorem}}{}
\@ifundefined{lemma}{\newtheorem{lemma}[theorem]{Lemma}}{}
\makeatother
\newcommand{\R}{\mathbb{R}}
\newcommand{\inter}[1]{#1^o}
\newcommand{\norme}[1]{\left\| #1 \right\|}
\newcommand{\scal}[2]{\left< #1,#2 \right>}
\newcommand{\tend}[2]{\underset{#1 \rightarrow #2}{\longrightarrow}}
\renewcommand{\geq}{\geqslant}
\title[Continuity of HYM connections with respect to metric variati]{Continuity of HYM connections with respect to metric variations}
\author{R\'emi Delloque}
\date{Source: arXiv:2403.16814v4 (CC BY 4.0)}
\begin{document}
\maketitle

\noindent\emph{Source and licence.} Continuity of HYM connections with respect to metric variations. Version arXiv:2403.16814v4 (CC BY 4.0), distributed under the Creative Commons Attribution 4.0 International licence, as recorded in the arXiv licence metadata for this exact version and shown on the abstract page https://arxiv.org/abs/2403.16814v4. Full rights notice: licenses/arxiv-2403.16814-CC-BY-4.0.txt.\par\smallskip

\noindent\textit{Editorial scope.} Counted unit: the Lemma whose source label is \texttt{LEM:Finitude locale} in the article (source0.tex lines 438--440, source label \texttt{LEM:Finitude locale}), delivered together with the article's own complete original proof of it (source0.tex lines 441--486, one \texttt{proof} environment of 46 lines). No mathematics is written, repaired or re-derived: statement and proof are the article's own text line for line. Required context retained uncounted: none. Every definition, display and label of those lines is retained unchanged. Omitted from this package and not redistributed: 1--437, 487--1440. Nothing inside a retained excerpt is omitted or altered: every retained body is delivered exactly as the article's lines are written, so no omission receipt of any kind accompanies this package. Citations: the retained excerpts carry no citation command at all, so no bibliography entry of the article is transcribed and no reference is redistributed. Reference adaptation: none was needed and none was made; every reference inside the delivered unit resolves inside it, so no unresolved reference or citation is produced. Printed slips kept literal: no printed word, symbol, hypothesis or number of the counted statement or of its proof was altered. Layout and class: the article's own class is \texttt{article} with packages including geometry, babel, amsmath, amsthm, amssymb, tikz-cd, dsfont, picture, float, graphicx, listings, hyperref, verbatim, todonotes; the wrapper is the lane's pinned \texttt{amsart} editorial wrapper at 10pt on A4, which restates the theorem-like environments the retained text needs and adds the article's own macro definitions that the retained text needs (\texttt{\textbackslash R}, \texttt{\textbackslash geq}, \texttt{\textbackslash inter}, \texttt{\textbackslash norme}, \texttt{\textbackslash scal}, \texttt{\textbackslash tend}), with extra wrapper packages (bm, bbm, dsfont, xspace, cancel, mathtools). The delivered numbering is the unit's own. Assets: no retained span contains a figure environment, a graphics-inclusion command, a PDF-page inclusion, a table or a TikZ picture; no image file is copied into this package, the package declares no asset dependency and no image file is redistributed with it. Licence: version arXiv:2403.16814v4 of this article is distributed under the Creative Commons Attribution 4.0 International licence, as recorded in the arXiv licence metadata for this exact version and shown on the abstract page https://arxiv.org/abs/2403.16814v4. A full rights notice is delivered at licenses/arxiv-2403.16814-CC-BY-4.0.txt. Characters: every retained excerpt is pure ASCII, so the delivered engine, XeLaTeX with its default Latin Modern text font, reports no missing character.
\begin{lemma}\label{LEM:Finitude locale}
    Let $V$ be a finite dimensional real vector space and $D \subset V$ a non-empty discrete subset. Let $U \subset V^\lor$ be an open subset of the dual of $V$ and such that for all $\varphi \in U$, $\varphi(D) \subset \R$ is bounded from above. Then, for all compact $K \subset U$, there is a finite set $F \subset D$ that may depend on $K$ such that for all $\varphi \in K$, $\sup(\varphi(D)) = \max(\varphi(F))$.
\end{lemma}

\begin{proof}\ \\
\vspace{1mm}
\noindent\textbf{Step 1 :} For all real number $a$, for all $K \subset U$ compact, the set $F_{K,a} = \bigcup_{\varphi \in K} \{\varphi \geq a\} \cap D$ is finite.
\vspace{1mm}

Indeed, $F_{K,a} \subset D$ is discrete so if it is infinite, then it is unbounded. In this case, there is a sequence $(v_m)$ of elements of $F_{K,a}$ such that $\norme{v_m} \rightarrow +\infty$ where $\norme{\cdot} = \sqrt{\scal{\cdot}{\cdot}}$ is any Euclidean norm on $V$. Up to extracting, $\frac{v_m}{\norme{v_m}} \rightarrow v_\infty \in V$ is in the unit sphere. For all fixed $m$, $v_m \in \bigcup_{\varphi \in K} \{\varphi \geq a\}$ so there is a $\varphi_m \in K$ such that $\varphi_m(v_m) \geq a$. Up to extracting again, $\varphi_m \rightarrow \varphi_\infty \in K$.

Let $\epsilon > 0$ and $\varphi = \varphi_\infty + \epsilon\scal{v_\infty}{\cdot}$. Since $U$ is open, for $\epsilon$ small enough, $\varphi \in U$. For all integer $m$,
\begin{align*}
    \varphi(v_m) & = \varphi_\infty(v_m) + \epsilon\scal{v_\infty}{v_m}\\
    & = \varphi_m(v_m) + \mathrm{O}(\norme{\varphi_m - \varphi_\infty}\norme{v_m}) + \epsilon\scal{\frac{v_m}{\norme{v_m}} + \mathrm{o}(1)}{v_m}\\
    & \geq a + \epsilon\norme{v_m} + \mathrm{o}(\norme{v_m})\\
    & \tend{m}{+\infty} +\infty.
\end{align*}
This contradicts the boundedness from above of $\varphi(D)$. It proves that $F_{K,a}$ is indeed finite. In particular, when $K = \{\varphi\}$ is a singleton and $a = \sup(\varphi(D)) - 1$, we obtain that for all $\varphi \in U$, the set,
$$
\{\varphi \geq \sup(\varphi(D)) - 1\} \cap D,
$$
is finite. In particular, it proves that the $\sup$ is reached (and it is reached finitely many times).

\vspace{1mm}
\noindent\textbf{Step 2 :} The map $\varphi \mapsto \max(\varphi(D))$ is continuous on $U$.
\vspace{1mm}

Let $\varphi_0 \in U$ and $K \subset U$ be a compact set such that $\varphi_0 \in \inter{K}$. Let $v \in D$ such that $\max(\varphi_0(D)) = \varphi_0(v)$ and $a = \max(\varphi_0(D)) - 1$. Clearly, $v \in F_{K,a}$ and when $\varphi$ is close enough to $\varphi_0$, then $\varphi \in K$ and $\varphi(v) \geq a$. Thus,
$$
\max(\varphi(D)) \geq \varphi(v) \geq a,
$$
so the $w \in D$ that verify $\max(\varphi(D)) = \varphi(w)$ belong to $F_{K,a}$. In other words, for all $\varphi$ close enough to $\varphi_0$, $\max(\varphi(D)) = \max(\varphi(F_{K,a}))$.

Since $F_{K,a}$ is finite, $\varphi \mapsto \max(\varphi(F_{K,a}))$ is continuous at $\varphi_0$, proving the wanted result.

\vspace{1mm}
\noindent\textbf{Step 3 :} Conclusion.
\vspace{1mm}

Let $K \subset U$ be an arbitrary compact set. By the continuity result of step 2,
$$
a = \min_{\varphi \in K}(\max(\varphi(D))),
$$
is well defined and by the result of step 1,
$$
F = F_{K,a} = \bigcup_{\varphi \in K} \{\varphi \geq a\} \cap D,
$$
is finite. For all $\varphi \in K$, if $v \in D$ is such that $\max(\varphi(D)) = \varphi(v)$, then $\varphi(v) \geq a$ so $v \in F$. We deduce that $\sup(\varphi(D))$ is reached by an element in the finite set $F$, which is the wanted result.
\end{proof}

\end{document}
